\section{Introduction}

Large language models (LLMs) trained on massive text collections have shown surprising emergent capabilities to generate text and perform zero- and few-shot learning~\cite{brown2020gpt3,J1WhitePaper,megatron2022,gopher2022,palm2022}. While in some cases the public can interact with these models through paid APIs, full model access is currently limited to only a few highly resourced labs.\footnote{Exceptions include work by EleutherAI, who released dense models up to 20B in size~\cite{eleutherai2022}, Salesforce \cite{Nijkamp2022ACP}, and Meta AI, who released dense models up to 13B and sparse models up to 1.1T~\cite{1TMoE2021}. There is also ongoing work from the BigScience workshop (\url{https://bigscience.huggingface.co/}), which aims to open source very large multilingual language models and datasets.} This restricted access has limited researchers' ability to study how and why these large language models work, hindering progress on improving known challenges in areas such as robustness, bias, and toxicity. 

In this technical report, we present Open Pre-trained Transformers (OPT), a suite of decoder-only pre-trained transformers ranging from 125M to 175B parameters, which we aim to fully and responsibly share with interested researchers. We train the OPT models to roughly match the performance and sizes of the GPT-3 class of models, while also applying the latest best practices in data collection and efficient training. Our aim in developing this suite of OPT models is to enable reproducible and responsible research at scale, and to bring more voices to the table in studying the impact of these LLMs.  Definitions of risk, harm, bias, and toxicity, etc., should be articulated by the collective research community as a whole, which is only possible when models are available for study.  

We are releasing all of our models between 125M and 66B parameters, and will provide full research access to OPT-175B upon request. Access will be granted to academic researchers; those affiliated with organizations in government, civil society, and academia; and those in industry research laboratories. 
We are also releasing both the logbook of our model creation as well as our codebase, {metaseq},\footnote{\url{https://github.com/facebookresearch/metaseq}} which enabled training \OPT{} on 992 80GB A100 GPUs, reaching ~147 TFLOP/s utilization per GPU.  From this implementation, and from using the latest generation of NVIDIA hardware, we are able to develop \OPT{} using only 1/7th the carbon footprint of \biggpt{.}  While this is a significant achievement, the energy cost of creating such a model is still nontrivial, and repeated efforts to replicate a model of this size will only amplify the growing compute footprint of these LLMs.

We believe the entire AI community — academic researchers, civil society, policymakers, and industry — must work together to develop clear guidelines around responsible AI in general and responsible LLMs in particular, given their centrality in many downstream language applications. A much broader segment of the AI community needs access to these models in order to conduct reproducible research and collectively drive the field forward. With the release of OPT-175B and smaller-scale baselines, we hope to increase the diversity of voices defining the ethical considerations of such technologies.
